\section{Conclusión}

El análisis documental confirma que la integración de paradigmas de programación, arquitecturas de software y metodologías ágiles constituye un eje fundamental en el desarrollo contemporáneo de sistemas. Según Introducción a la Programación Orientada a Objetos \cite{poo}, la POO facilita modularidad y adaptabilidad en el código, mientras que el patrón Modelo–Vista–Controlador (MVC) aporta organización y escalabilidad en las aplicaciones \cite{mvc}. Por su parte, Node.js vs Spring Boot en el desarrollo backend de aplicaciones web \cite{nodejs_springboot} evidencia que la elección de frameworks debe responder a las necesidades específicas del proyecto, equilibrando rapidez y flexibilidad frente a seguridad y robustez.

Asimismo, la Metodología Scrum \cite{scrum} refuerza la comunicación del equipo, la planificación y la entrega incremental de valor, y el libro Fundamentos sobre la gestión de bases de datos \cite{databases} resalta la importancia de modelos sólidos para garantizar la consistencia e integridad de la información.

En conjunto, la combinación de la Programación Orientada a Objetos, la arquitectura MVC, frameworks modernos como Node.js y Spring Boot, metodologías ágiles y bases de datos robustas permite construir sistemas más escalables, seguros y alineados con las demandas actuales de la industria del software.